% !TeX root = ../ADCS & GNC Documentation.tex
\chapter{Concept of Operations}

\section{Overview}

The ADCS system is designed to operate with minimal user interaction. From an operational perspective, the system requires only high-level mission selection and execution parameters, while all attitude control, guidance, and communication alignment are handled autonomously onboard.

\section{Mission Definition}

Operation of the system consists of selecting:

\begin{itemize}
	\item A mission type (e.g. nadir pointing, or max-drag/ram-facing mode)
	\item A duration for which the mission should be executed
\end{itemize}

Once these parameters are defined, the ADCS system executes the mission autonomously for the specified duration.

\section{Attitude Modes}

Two primary modes are used:

\begin{itemize}
	\item \textbf{Nadir Pointing:} The spacecraft maintains its +z axis aligned toward the center of the Earth. This is suitable for maintaining a consistent ground-facing orientation.
	
	\item \textbf{Ram-Facing Orientation:} The spacecraft aligns it's +x axis with the velocity vector (ram direction). This results in the radar-occlusion GPS antenna facing anti-ram.
\end{itemize}

For circular or near-circular orbits, the difference between these modes is minimal, as the velocity vector is nearly perpendicular to the radial (nadir) direction.

In elliptical orbits, however, the velocity vector varies significantly along the orbit, particularly near perigee and apogee. In these cases, ram-facing and nadir-pointing orientations diverge.

\section{Autonomous Overpass Detection and Downlink}

Ground station communication is handled autonomously using pre-defined ground station locations.

After mission completion, if immediate downlink is requested:

\begin{itemize}
	\item The spacecraft monitors its orbit to determine upcoming ground station overpasses
	\item The next available overpass window is identified
	\item The spacecraft transitions into the appropriate communication attitude mode (target tracking) once it nears the overpass window
	\item Data transmission begins once the ground station is within range
\end{itemize}

This process does not require real-time user intervention and is driven by onboard prediction of orbital position relative to known ground station locations.

\section{Operational Summary}

From a user perspective, operation of the ADCS system consists of:

\begin{enumerate}
	\item Selecting a mission mode
	\item Specifying the mission duration
	\item Indicating whether immediate data should be downloaded as soon as possible after mission completion
\end{enumerate}

All subsequent attitude control, guidance, and communication alignment are handled autonomously by the system.